% =====================================================================
\chapter{Memoria, segundo orden y método adjunto}
\label{tema:5}
% =====================================================================

El modo inverso compra su eficiencia en operaciones con memoria: necesita los valores de la pasada hacia delante. Este tema cuantifica esa memoria y estudia el recálculo por puntos de control, que la reduce a cambio de cómputo; obtiene productos hessiana-vector combinando los dos modos, y muestra que los adjuntos son multiplicadores de Lagrange, lo que conecta la retropropagación con el método adjunto del control óptimo y de la física. Se siguen los capítulos~4 y~12 de~\cite{gw}, el artículo de Chen \emph{et al.}~\cite{chen}, el de Pearlmutter~\cite{pearlmutter}, la §3.3 de~\cite{baydin} y la tarea~2 de CS336~\cite{cs336a2}.

\section{La memoria del modo inverso}

\begin{proposicion}[Memoria de una cadena de capas]\label{prop:memoria}
Sea $f=f_L\circ\dots\circ f_1$ una red de $L$ capas cuyas activaciones intermedias $h_k=f_k(h_{k-1})$ ocupan $a$ unidades de memoria cada una, y supóngase que el VJP de la capa $k$ necesita $h_{k-1}$. Evaluar $f$ requiere memoria $O(a)$, mientras que el modo inverso de la \cref{fig:algoritmo} requiere guardar $h_0,\dots,h_{L-1}$, es decir, memoria $La$.
\end{proposicion}

En un LLM, $a$ es proporcional al número de tokens del lote por la dimensión del modelo, y la memoria de las activaciones supera a menudo la de los pesos (módulo~17).

\section{Recálculo por puntos de control}

La idea del recálculo (\emph{gradient checkpointing} o \emph{rematerialization}) consiste en guardar solo algunas activaciones, los \emph{puntos de control}, y recalcular las demás cuando la pasada hacia atrás las necesita~\cite{chen}, \cite[cap.~12]{gw}.

\begin{proposicion}[Un nivel de puntos de control]\label{prop:checkpoint}
En la situación de la \cref{prop:memoria}, divídanse las $L$ capas en $L/k$ tramos de $k$ capas consecutivas y guárdese durante la pasada hacia delante solo la entrada de cada tramo. En la pasada hacia atrás, al llegar a cada tramo se repite su pasada hacia delante desde el punto de control, guardando sus $k$ activaciones, y se retropropaga por él. Entonces:
\begin{enumerate}
  \item la memoria de activaciones es $M(k)=(L/k+k)\,a$, mínima para $k=\sqrt L$, con valor $2\sqrt L\,a$;
  \item el cómputo adicional es una pasada hacia delante más.
\end{enumerate}
\end{proposicion}

Con las cuentas de la \cref{prop:flops}, una pasada hacia delante extra cuesta $2N$ FLOPs por token, y el paso de entrenamiento pasa de $6N$ a $8N$: un 33\,\% más de cómputo a cambio de reducir la memoria de $L$ a $2\sqrt L$ activaciones. Aplicando la idea de forma recursiva, la memoria baja a $O(\log L)$ con un cómputo de $O(L\log L)$ evaluaciones de capa~\cite{chen}; Griewank y Walther estudian los esquemas óptimos de este tipo~\cite[cap.~12]{gw}.

\addtocounter{definicion}{1}%

\section{Productos hessiana-vector}

Las derivadas de segundo orden se obtienen aplicando la AD a un programa que, a su vez, calcula derivadas.

\begin{proposicion}[Directo sobre inverso]\label{prop:hvp}
Sea $f\colon\R^{n}\to\R$ de clase $C^{2}$ dada por una traza cuyas operaciones elementales son de clase $C^{2}$, y sean $x,v\in\R^{n}$. Si se ejecuta el modo inverso (que calcula $\nabla f$) en aritmética dual, con entradas $x+v\varepsilon$, la parte dual del resultado es el producto hessiana-vector $H_f(x)\,v$. Su coste es a lo sumo una constante, independiente de $n$, por $\ops(f)$.
\end{proposicion}

\addtocounter{definicion}{1}%

La técnica es la de Pearlmutter~\cite{pearlmutter} y permite usar información de curvatura sin formar el hessiano, que para $n=10^{9}$ tendría $10^{18}$ entradas. El hessiano completo requiere $n$ productos hessiana-vector, uno por columna, y cuesta del orden de $n\cdot\ops(f)$, frente al $O(\ops(f))$ del gradiente.

\section{Los adjuntos como multiplicadores de Lagrange}

\begin{teorema}[Método adjunto]\label{teo:lagrange}
Sea una traza de $f\colon\R^{n}\to\R$ y considérese la lagrangiana
\[
\mathcal L(x,v_1,\dots,v_\ell,\lambda_1,\dots,\lambda_\ell)=v_\ell-\sum_{i=1}^{\ell}\lambda_i\Bigl(v_i-\varphi_i\bigl((v_j)_{j\pre i}\bigr)\Bigr),
\]
donde $v_{k-n}=x_k$. Si los $v_i$ toman los valores de la traza y los multiplicadores se eligen de modo que $\partial\mathcal L/\partial v_i=0$ para $i=1,\dots,\ell$, entonces $\lambda_i=\bar v_i$ para todo $i$ y $\nabla f(x)=\nabla_x\mathcal L$.
\end{teorema}

\addtocounter{definicion}{1}%

El ejemplo es la retropropagación a través del tiempo de una red recurrente: un mismo parámetro actúa en cada paso y su gradiente es la suma de las contribuciones de todos ellos. En control óptimo, en diseño aerodinámico y en la asimilación de datos meteorológicos (el método 4D-Var), el mismo esquema se aplica a ecuaciones diferenciales, integrando una ecuación adjunta hacia atrás en el tiempo para obtener el gradiente de un coste respecto a millones de condiciones iniciales o parámetros.

\begin{observacion}[Historia]
Linnainmaa describió el modo inverso en 1970 como herramienta para analizar errores de redondeo acumulados~\cite{linnainmaa}; Rumelhart, Hinton y Williams lo popularizaron en 1986 para entrenar redes neuronales~\cite{rumelhart}. Baydin \emph{et al.} reconstruyen esta historia y sus redescubrimientos~\cite[§3.3]{baydin}.
\end{observacion}

